\section{The integral Jacquet--Langlands defect}\label{sec:stage-06}
% Stage 06 (integral-defect). Drafter: Opus stage drafter. Source: seed §9 (and the last paragraph of §10, O5, C1, C3, C4(i),
% T4, T5).
\providecommand{\BG}{\mathrm{BG}}          % bipartite Brandt graph of (U,q); same definition as stages 02/03

\subsection{Goals}
This stage makes seed \S9 precise. It works on the $\GL_2$ side, with the definite algebra $B=B_q$ entering only through the
$\Lambda$-adic Hecke JL $\T'\cong\T^{\ord}_{B}$ and Hsieh's description of $\eta^B_F$. It does not use $\HypX$. The goals are:
\begin{enumerate}[label=(G\arabic*)]
\item \emph{Commutative algebra with all hypotheses explicit.} Lemmas S, G, D and Thm~E (the combinatorial half
  $\eta_F=\cc_q(F)\cdot\eta^B_F$). Thm~E needs \emph{two} Gorenstein hypotheses: (Gor$_R$) on the full algebra $R$, which is used
  through Lemmas S and D, and (Gor$_{R'}$) on the $q$-new quotient. Without (Gor$_{R'}$) the identity still holds at height-one primes.
  We also record a weaker sufficient condition (Thm~\ref{thm:s06-E}(ii)).
\item \emph{Galois half.} $\tq_q(F)\subseteq\cc_q(F)$ (Prop.~\ref{prop:s06-galois}), with the inputs it needs stated precisely.
\item \emph{Defect formula.} $(h_F)=\cc_q(F)=\tq_q(F)$ (Thm~\ref{thm:s06-defect}). Step (4) of the seed has two proofs.
  [KO, Thm~1.1] as a black box excludes $q\equiv-1\bmod p$ when $\bar\rho$ is unramified at $q$; this gives the theorem under
  (H0$^-$) $q\not\equiv-1\bmod p$. The seed's direct $R=\T$ argument (Prop.~\ref{prop:s06-levM}) does not need it; it gives the
  theorem under (Lev$_M$), which holds for $M=1$ (Remark~\ref{rem:s06-minus1}). The seed's weight-two alternative in (P) also needs
  one more condition, (P2).
  \provisional{F3/D2: two proofs of step (4), theorem under (H0$^-$) or (Lev$_M$); before: ``the seed's step (4) is corrected: Kim--Ota
  excludes $q\equiv-1$, so we add (H0$^-$)''.}
\item \emph{All weights.} Cor.~\ref{cor:s06-allweights} extends Kim--Ota's fixed-weight identity to all arithmetic weights $k'\ge3$
  of the branch, including $k'\ge p$.
\item \emph{Betti incarnation.} The $\GL_2$ level-raising module $\coker(\Psi^*\Psi)_\m$ of stage 03 is identified with
  $(\T(M)_\m/\det\Psi^*\Psi)^{\oplus2}$ when $q\not\equiv-1\bmod p$. It computes $\cc_q$ when $\rk\T'_\m=1$ (Prop.~\ref{prop:s06-betti}).
  This is the basis of T4.
\item \emph{Beyond $D=q$.} The open cases (O5) and the conjectural several-primes formula (seed \S10, last paragraph).
\end{enumerate}

\subsection{Definitions and statements}

\subsubsection{Setting and standing hypotheses}\label{ss:s06-setting}
Here $D=q$, so $B=B_q$ is ramified at $\{q,\infty\}$. $p\ge5$, $q\ne p$, $(M,pq)=1$, $N:=qM$. $\Lam=\cOE[[1+p\Zp]]$, with
uniformiser $\varpi$ of $\cOE$. On the $\GL_2$ side, $\Lam$ acts through the \emph{classical} diamond operators, and
$P_k=\ker([u]\mapsto u^{k-2})$ is the weight-$k$ point (seed \S2.6). On the definite side, stage 04 lets $\Lam$ act by
\emph{adelic} diamonds. These are inverse to the classical ones (stage 01). In those conventions the weight-$k$ point is
$P^{\rm ad}_k=\ker([u]\mapsto u^{2-k})$ (\texttt{if:s04-Mord}, stage 04 D1). The $\Lam$-adic JL $\T'\cong\T^{\ord}_B$ is
$\Lam$-linear only if the two actions are matched by $[u]\mapsto[u]^{-1}$. We always use the $\GL_2$ normalisation and transport it.

\begin{hypothesis}[standing hypotheses; seed \S9.1, corrected]\label{hyp:s06-H}
\leavevmode
\begin{itemize}
\item[(H0)] $p\ge5$, $q\ne p$, $(M,pq)=1$, $N=qM$, $p\nmid\varphi(N)$. In particular $q\not\equiv1\bmod p$.
\item[(H0$^-$)] (\emph{new}; needed for the fixed-weight input Cor.~\ref{cor:s06-fixed} and for Prop.~\ref{prop:s06-betti};
  Thm~\ref{thm:s06-defect} needs (H0$^-$) \emph{or} (Lev$_M$); the inclusion of Prop.~\ref{prop:s06-galois} uses neither)
  $q\not\equiv-1\bmod p$, or
  $\bar\rho_\m$ is ramified at $q$. Together with (H0): $p\nmid q^2-1$, or $\bar\rho_\m$ is ramified at $q$.
\item[(Lev$_M$)] (\emph{new}; alternative to (H0$^-$) in Thm~\ref{thm:s06-defect}) At the point $P$ of (P), of weight $k=k_P$:
  the level-$M$ Hecke algebra $\T_k(M)_\m$ is the Hecke algebra of Kim--Ota's refined $R=\T$ theorem [KO, Thm~3.14]
  \lit{kim-ota\_2019\_*}{Thm~3.14, p.~12} for $\Sigma^-=\{\ell\,\|\,M:\bar\rho\text{ unramified at }\ell\}$, $\Sigma^+=\emptyset$,
  and the hypotheses of that theorem hold (Assumption~2.1 of [KO] and $P_{\bar\rho}\subseteq\Sigma$). It holds for $M=1$: then
  nothing is needed beyond Assumption~2.1, which is (H2) plus the existence of $f$.
  \provisional{F3/D2: hypothesis (Lev$_M$) added and (H0$^-$) made one of two alternatives; before: (H0$^-$) alone was ``needed
  for the equality in Thm defect''.}
\item[(H1)] $\T$ is Hida's full ordinary cuspidal Hecke algebra of tame level $\Gamma_0(N)$ over $\Lam$, in a fixed weight component
  $k_0\bmod(p-1)$. It is generated by all $T_\ell$ ($\ell\nmid Np$), $U_\ell$ ($\ell\mid Np$, in particular $U_q$ and $U_p$), and it
  is localised at a non-Eisenstein maximal ideal $\m$. $F$ is a primitive ordinary $q$-new branch, with
  $\lambda_F:\T\to I_F$, and $I_F$ is a normal local domain finite over $\Lam$. Then $I_F$ is $\Lam$-free: it is normal of dimension $2$, hence
  Cohen--Macaulay (Serre's criterion), hence free over the regular local ring $\Lam$ (Auslander--Buchsbaum). So the seed's
  ``hence $\Lam$-free'' is correct; it is used only through flatness of $R=\T\otimes_\Lam I_F$ over $I_F$.
\item[(H2)] $\bar\rho=\bar\rho_\m$ is absolutely irreducible on $G_{\Q(\sqrt{p^*})}$ and $p$-distinguished.
\item[(Im)] $p\mid\#\bar\rho(G_\Q)$. \emph{None of the cited theorems below uses (Im).} [KO, Thm~1.1] needs only (H2) (its
  condition (2) = Assumption 2.1(TW)), and [KO, Rem.~2.10] says that the [DFG] and [Dim09] image conditions are interchangeable over
  $\Q$. We keep (Im) as in the seed, flagged as probably superfluous (Risk R4).
\item[(P)] Either (P1) $k_0\not\equiv2\bmod(p-1)$, and then $F$ has an arithmetic point $P$ of weight $k_P\in[3,p-1]$ with trivial
  character (a \emph{Kim--Ota point}). Or (P2) $k_0\equiv2$, $F$ has a weight-two point $P$ which is the ordinary
  $p$-stabilisation of a newform $f$ of level $N$, \emph{and} (\emph{new}) the weight-two ordinary Hecke algebra
  $\T^{\ord}_2(\Gamma_0(Np))_\m$ has no $p$-new eigenforms. The extra condition makes $\T/P\T\cong\T_2(\Gamma_0(N))_\m$,
  which is the algebra of [KO]. Without it the control step (3) of the proof compares with an algebra that [KO] does not treat.
\end{itemize}
\end{hypothesis}

\begin{remark}[relation to Hsieh's residual condition]\label{rem:s06-CR}
Hsieh's (CR,$\Sigma$) is \lit{hsieh\_2017\_*}{\S1.4, p.~4}: $\bar\rho_F$ is absolutely irreducible and $p$-distinguished, and
it is ramified at every $\ell\in\Sigma$ with $\ell\equiv1\bmod p$. The congruence is $+1$ only, not $\pm1$ (stage 04 point (4)). With
$\Sigma=\Sigma^-=\{q\}$, (H2)$+$(H0) imply (CR,$\{q\}$), because absolute irreducibility on $G_{\Q(\sqrt{p^*})}$ implies it on
$G_\Q$ and (H0) gives $q\not\equiv1$. Hence the \emph{combinatorial half needs only (H0), (H1), (H2)}. The congruence
$q\equiv-1$ enters only through the fixed-weight Galois input (Remark~\ref{rem:s06-minus1}).
\end{remark}

\subsubsection{Hecke algebras, congruence ideals, defect}\label{ss:s06-rings}
\begin{definition}[seed \S9.1]\label{def:s06-rings}
$\T'$ and $\T''$ are the images of $\T$ in the endomorphisms of the $q$-new and of the $q$-old ordinary $\Lam$-adic cusp forms
(localised at $\m$). Put $\fa:=\ker(\T\to\T')$ and $\fb:=\ker(\T\to\T'')$. Put $R:=\T\otimes_\Lam I_F$, $R':=\T'\otimes_\Lam I_F$,
$R'':=\T''\otimes_\Lam I_F$, and keep the names $\fa,\fb$ for $\fa R,\fb R$. $\lambda:R\to I_F$ is the $I_F$-linear extension of
$\lambda_F$. Since $F$ is $q$-new, $\lambda$ factors as $\lambda=\lambda'\circ(R\to R')$. For a finite $A$-algebra $S$ and an
$A$-algebra map $\mu:S\to A$ to a domain, $\eta(S,\mu):=\mu(\Ann_S\ker\mu)$ (seed \S2.6). Put
\[\eta_F:=\eta(R,\lambda),\qquad \eta^B_F:=\eta(R',\lambda'),\qquad (h_F):=\eta_F\,(\eta^B_F)^{-1}\ \text{ when both are principal.}\]
\status{definition}
\end{definition}
The seed writes $h_F=\eta_F/\eta^B_F$ as an element. This makes sense only when both ideals are principal, which needs (Gor$_R$) and
(Gor$_{R'}$) (Lemma~\ref{lem:s06-S}). Under them, $h_F\in I_F$ is well defined up to $I_F^\times$ (Thm~\ref{thm:s06-E}).

\begin{lemma}[old/new decomposition]\label{lem:s06-decomp}
Assume (H1). (a) $\T\otimes_\Lam K=(\T'\otimes K)\times(\T''\otimes K)$ with $K=\mathrm{Frac}\,\Lam$, $\fa\cap\fb=0$, and
$\T/\fa$, $\T/\fb$ are $\Lam$-torsion-free. (b) Let $\T(M)_\m$ be Hida's ordinary algebra of tame level $M$ localised at the
pull-back of $\m$, and $S_q\in\Lam^\times$ the element with $P_k(S_q)=q^{k-2}$ (trivial character). If $\bar\rho_\m$ is ramified at
$q$, then $\T''=0$. Otherwise $X^2-T_qX+qS_q$ has two roots in $\T(M)_\m$ that are distinct modulo the maximal ideal: they are
$\equiv U_q$ and $\equiv qU_q$ mod $\m$, and $q\not\equiv1$ by (H0). Let $\alpha_q$ be the root congruent to $U_q$. Then
$\T''\cong\T(M)_\m$ via $U_q\mapsto\alpha_q$. In particular $\T''$ is $\Lam$-free, and $T_q=U_q+qS_qU_q^{-1}$ lies in the image
of $\T$. The same holds at fixed weight $k$: $\T''_k\cong\T_k(M)_\m$, $U_q\mapsto\alpha_q$.
\provisional{F2: the old quotient is the localisation $\T(M)_\m$, not $\T(M)_\m[X]/(X^2-T_qX+qS_q)$ (which has twice the rank);
before: ``$\T''\cong\T(M)_\m[X]/(X^2-T_qX+qS_q)$, $X\mapsto U_q$; in particular $\T''$ is $\Lam$-free''.} (c) The same statements hold after
$\otimes_\Lam I_F$ for $R,R',R''$. \sttoprove{} (standard)
\end{lemma}
\noindent\emph{Inputs.} (a) Strong multiplicity one separates $q$-new and $q$-old eigensystems at every arithmetic point. Hence the
generalised eigenspaces of $\T\otimes K$ are either new or old, and $\T$ is torsion-free because it acts faithfully on a $\Lam$-free
module. (b) The old space is two copies of the level-$M$ space with $U_q$ acting by the companion matrix of $X^2-T_qX+qS_q$. Old forms
are unramified at $q$, so there are none at $\m$ if $\bar\rho_\m$ is ramified at $q$. Otherwise the polynomial factors by Hensel
($\T(M)_\m$ is complete local and the two roots differ mod $\m$), so $\T(M)_\m[X]/(X^2-T_qX+qS_q)\cong\T(M)_\m\times\T(M)_\m$, and
localising at $\m$ (where $U_q\equiv\alpha_q$) keeps one factor. $\Lam$-freeness of $\T(M)_\m$ is [Hi86b, Thm~3.1]
\lit{hida\_1986\_iwasawa-modules-attached-to-congruences-of-cusp-forms}{Thm 3.1, p.~250} ($p\ge5$). $T_q=\alpha_q+qS_q\alpha_q^{-1}$ is
then the image of $U_q+qS_qU_q^{-1}$. \provisional{F2: Hensel/localisation argument replaces the Chebotarev argument for ``$T_q$ in
the image''; before: the subtle point was surjectivity onto $\T(M)_\m[X]/(\cdots)$, proved by topological generation of $\T(M)_\m$
by the $T_\ell$, $\ell\ne q$ (Chebotarev).} [Hi86b, Cor.~3.3, p.~250] gives $\Lam$-freeness of the old \emph{ideal}, not of $\T''$. Hida's old
ideal is the annihilator of the forms old at \emph{any} divisor $t>1$ of the tame level [Hi86b, (3.4a--d), p.~250], so it is contained in
our $\fb$ ($q$-old only) and equals it only when no form in the $\m$-part is old at a prime of $M$ (e.g.\ $M=1$).
(c) Flat base change, using (a), (b).

\begin{definition}[combinatorial Tamagawa ideal; seed Def.~ctam]\label{def:s06-ctam}
$\cc:=\mathrm{im}(\fb\to R')=(\fa+\fb)/\fa$, so $R'/\cc\cong R/(\fa+\fb)$ is the $\Lam$-adic old/new congruence module at $q$
(after $\otimes I_F$). Put $\cc_q(F):=\lambda'(\cc)=\lambda(\fb)\subseteq I_F$. At an arithmetic point $P$ with $q$-new
specialisation $f$, the \emph{fixed-weight} version is $\cc_q(f):=\lambda_f(\fb_k)$, where $\fb_k=\ker(\T_k(N)_\m\to\T''_k)$ is taken
in the weight-$k$ Hecke algebra. \status{definition}
\end{definition}
\noindent This is Diamond's $\Lam$-adic old/new congruence module [Dia89]. Diamond computes its specialisation along the \emph{old}
branch: [Dia89, Thm~6B, p.~80]
\lit{diamond\_1989\_congruence-modules-associated-to-lambda-adic-forms}{Thm 6A--6C, p.~80} \stcited{} gives
$\mathrm{div}(\mathfrak b)=-\mathrm{div}(a_q-\xi(q+1))$ for the sign-$\xi$ part. Here we use the \emph{new} branch, where the same
module is read through $\lambda'$ instead.

\begin{definition}[Tamagawa ideal; seed \S9.1]\label{def:s06-tam}
Let $H$ be Hida's ordinary $\Lam$-adic étale cohomology of tame level $N$ localised at $\m$. It is free of rank two over $\T$
under (H2) (cited below). Put $\mathcal T_F:=H\otimes_{\T,\lambda_F}I_F$, with $\rho_F:G_\Q\to\GL_2(I_F)$. Define
$\tq_q(F)\subseteq I_F$ as the ideal generated by the entries of $\rho_F(\sigma)-1$, $\sigma\in I_q$. (We write $\mathcal T_F$
rather than the seed's $\cT_F$, since $\cT$ is the tree.) \status{definition}
\end{definition}

\begin{lemma}\label{lem:s06-tam}
(a) $\rho_F|_{I_q}$ is unipotent and factors through the $\Zp(1)$-quotient of tame inertia. So $\tq_q(F)$ is generated by the
entries of $\rho_F(\tau)-1$ for one generator $\tau$. (b) At an arithmetic point $P$ with $q$-new $f=F_P$,
$P(\tq_q(F))\,\cO_P=(\varpi_P^{t_f(q)})$, where $\cO_P=I_F/P$ and $t_f(q)$ is the Tamagawa exponent of
[KO, Def.~2.12] \lit{kim-ota\_2019\_*}{Def.~2.12, p.~6} (= [PW, Def.~3.3]
\lit{pollack-weston\_2006\_*}{Def.~3.3, p.~9}). (c) $t_f(q)<\infty$ because $f$ is $q$-new. \sttoprove{} (standard)
\end{lemma}
\noindent (a) holds at every arithmetic point by local--global compatibility (Steinberg at $q$), and the arithmetic points are dense.
(b) $\mathcal T_F/P$ is a $G_\Q$-stable $\cO_P$-lattice in $V_f$, which is unique up to scaling by (H2). Ramification indices of
$\cO_P$ over $\cO_{f,\lambda}$ rescale both sides equally.

\subsubsection{Commutative algebra}\label{ss:s06-ca}
Let $A$ be a Noetherian local domain with fraction field $K$. A finite free $A$-algebra $S$ is \emph{$A$-Gorenstein} if
$S^\vee:=\Hom_A(S,A)$ is free of rank one over $S$. A generator $\ell\in S^\vee$ is called a \emph{trace}.

\begin{lemma}[base change of congruence ideals; seed Lemma S]\label{lem:s06-S}
Let $S$ be finite free and $A$-Gorenstein with trace $\ell$, and let $\mu:S\to A$ be an $A$-algebra map. Then
$S[\mu]:=\Ann_S\ker\mu=A\theta$ is free of rank one, with $\ell(\theta x)=\mu(x)$. Hence $\eta(S,\mu)=(\mu(\theta))$ is principal.
For $\varphi:A\to A'$ local, $S\otimes A'$ is $A'$-Gorenstein and $\eta(S\otimes A',\mu\otimes A')=\varphi(\eta(S,\mu))A'$. \stproved{}
(seed; proof: $S[\mu]\cong\Hom_S(S/\ker\mu,S^\vee)=\Hom_A(A,A)$, and all of this commutes with base change since $S$ is free.)
\end{lemma}

\begin{lemma}[Gorenstein from one fibre; seed Lemma G]\label{lem:s06-G}
Let $S$ be finite free over $A$ and $\mathfrak p\subseteq\m_A$. If $S/\mathfrak pS$ is $A/\mathfrak p$-Gorenstein, then $S$ is
$A$-Gorenstein. \stproved{} ($S^\vee/\mathfrak pS^\vee=(S/\mathfrak pS)^\vee$ is cyclic, so $S^\vee$ is cyclic by Nakayama over
the semilocal $S$; then $S\twoheadrightarrow S^\vee$ is a surjection of free $A$-modules of equal rank.)
\end{lemma}
\noindent The seed's ``$S$ local'' is not needed; $S$ finite over local $A$ suffices. This lemma is how Hsieh's $\Lam$-adic
Gorenstein statements follow from the weight-two ones of [Wi95, Cor.~2 to Thm~2.1, p.~482]
\lit{wiles\_1995\_modular-elliptic-curves-and-fermats-last-theorem}{Thm 2.1 and Cor.~2, p.~482} plus control.

\begin{lemma}[new/old duality; seed Lemma D, corrected]\label{lem:s06-D}
Let $S$ be finite free over $A$, with ideals $\fa,\fb$ such that $S/\fa$, $S/\fb$ are $A$-torsion-free, $\fa\cap\fb=0$ and
$S_K=(S/\fa)_K\times(S/\fb)_K$. Then:
\begin{enumerate}[label=(\alph*)]
\item $\Ann_S(\fa)=\fb$ and $\Ann_S(\fb)=\fa$. (No Gorenstein hypothesis is needed.)
\item If $S$ is $A$-Gorenstein with trace $\ell$, then $\beta\mapsto\ell(\beta\,\cdot)$ gives $S$-isomorphisms
  $\fb\cong(S/\fa)^\vee$ and $\fa\cong(S/\fb)^\vee$.
\item If moreover $S/\fb$ is $A$-free, then $\fb$ is a direct summand of $S$ as an $A$-module, hence $A$-free, and
  $s\mapsto\ell(s\,\cdot)|_\fb$ induces $S/\fa\cong\fb^\vee$. In particular $S/\fa$ is $A$-free.
\end{enumerate}
All of this commutes with base change $A\to A'$ that preserves the hypotheses. \stproved{} ((a): $\Ann(\fa)=S\cap(S'_K\times0)=\fb$.
(c): $\ell(s\fb)=0\Rightarrow s\fb=0$, so the kernel is $\Ann\fb=\fa$; surjectivity holds because $\fb$ is a summand.)
\end{lemma}
\noindent The seed states (c) without saying \emph{why} $S/\fa$ is free. Freeness of the dual $(S/\fa)^\vee\cong\fb$ alone does not
imply it when $\dim A\ge2$. The argument above uses the summand property, i.e.\ freeness of $S/\fb$ (Lemma~\ref{lem:s06-decomp}(b)).

\begin{theorem}[the combinatorial half; seed Thm~E, hypotheses explicit]\label{thm:s06-E}
Let $S,\fa,\fb$ be as in Lemma~\ref{lem:s06-D}. Put $S'=S/\fa$, and let $\lambda:S\to A$ factor through $\lambda':S'\to A$. Assume
\begin{itemize}
\item[(Gor$_S$)] $S$ is $A$-Gorenstein.
\end{itemize}
Let $\theta$ generate $S[\lambda]$ (Lemma~\ref{lem:s06-S}), and let $x\in S'$ be its image. Then:
\begin{enumerate}[label=(\roman*)]
\item $\theta\in\fb$, $x\in S'[\lambda']$, $\eta(S,\lambda)=(\lambda'(x))$, and $\cc_q:=\lambda(\fb)=\{\psi(x):\psi\in S'^\vee\}$ (the
  \emph{order ideal} of $x$).
\item If $S'[\lambda']$ is free of rank one with a generator $\theta'$ that is unimodular in $S'$ (i.e.\ $S'[\lambda']$ is an
  $A$-direct summand of $S'$), then $\eta(S,\lambda)=\cc_q\cdot\eta(S',\lambda')$.
\item This holds in particular under (Gor$_{S'}$): $S'$ is finite free and $A$-Gorenstein. It also holds whenever $A$ is a DVR
  (\emph{no} hypothesis on $S'$). For $A$ normal, the identity therefore holds after localisation at every height-one prime.
\item In general $\cc=\mathrm{im}(\fb\to S')\cong\fb\cong S'^\vee$. So $\cc$ is a principal ideal of $S'$ if and only if
  (Gor$_{S'}$) holds.
\end{enumerate}
With $S=R$, $A=I_F$: under (Gor$_R$) and (Gor$_{R'}$), $\eta_F=\cc_q(F)\cdot\eta^B_F$, i.e.\ $(h_F)=\cc_q(F)$. \stproved{}
\end{theorem}
\noindent\emph{Proof} (complete; it corrects the seed's sketch). $\fa\subseteq\ker\lambda$ gives $\theta\in\Ann(\fa)=\fb$.
For $s\in S$, $\theta s=\lambda(s)\theta$, so $\ell(\theta\,\cdot)=\ell(\theta)\lambda'$ on $S'$. Lemma~\ref{lem:s06-D}(b) restricts to
$\fb[\lambda]=S[\lambda]\xrightarrow\sim S'^\vee[\lambda']=A\lambda'$, so $\ell(\theta)\in A^\times$; normalise $\ell(\theta)=1$.
For $\beta\in\fb$, $\lambda(\beta)=\ell(\beta\theta)=\psi_\beta(x)$, and $\psi_\beta$ runs through $S'^\vee$. This proves (i). For (ii),
write $x=h\theta'$. Then $\lambda'(x)=h\lambda'(\theta')$ and $\{\psi(x)\}=h\{\psi(\theta')\}=hA$. For (iii): under (Gor$_{S'}$),
$\ell'(\theta')\in A^\times$ by the same argument, so $\theta'$ is unimodular. Over a DVR, $S'[\lambda']=S'\cap(K\times0)$ is
saturated, hence a summand. (iv) is Lemma~\ref{lem:s06-D}(b) together with $\fa\cap\fb=0$. \qed

\noindent\emph{Where each Gorenstein hypothesis is used.} (Gor$_R$) gives $\theta$ and the principality of $\eta_F$
(Lemma~\ref{lem:s06-S}), and the isomorphism $\fb\cong R'^\vee$ (Lemma~\ref{lem:s06-D}(b)). It is \emph{indispensable}: without it,
$\eta_F$ need not be principal and $\fb\not\cong R'^\vee$. (Gor$_{R'}$) is used only in (ii)$\Leftarrow$(iii), to make the
generator of $R'[\lambda']$ unimodular. The seed's phrase ``in the setting of Lemma D'' hides (Gor$_R$).
\emph{Sources.} (Gor$_R$) is [Hs, \S3.3, p.~18] \lit{hsieh\_2017\_*}{\S3.3, p.~18; \S7.2, p.~68}
\stcited{} under (CR), via [Wi95, Cor.~2 p.~482] and [MW86, \S9 Prop.~2] (MW86 is not in \texttt{literature/}). (Gor$_{R'}$) is
\texttt{if:s04-Mord} ([Hs, Thm~4.5 and proof, p.~34]: $\T^D(N,\chi,\mathbb I)_\m$ is Gorenstein) transported by the JL isomorphism
$\T'\cong\T^{\ord}_B$ ([Hs, Rem.~7.8, p.~71]: $\T^D$ is the $N^-$-new quotient). Both are then base-changed to $I_F$ by
Lemma~\ref{lem:s06-S}. Consistency: Thm~E gives $\eta_F\subseteq\eta^B_F$, which is the inclusion $(\eta^B_F)\supseteq(\eta_F)$
of [Hs, Rem.~7.8] (\texttt{if:s04-etaB}).

\subsubsection{The Galois half}
\begin{proposition}[Galois comparison; seed Prop.~galois]\label{prop:s06-galois}
Assume (H0), (H1), (H2). Then $\tq_q(F)\subseteq\cc_q(F)$. \stpmc{}
\end{proposition}
\noindent\emph{Cited inputs.} (1) $H$ is free of rank two over $\T$, and $H(M)$ is free of rank two over $\T(M)_\m$ (tame level $M$):
[Hs, \S7.2, p.~68] \lit{hsieh\_2017\_*}{\S7.2, proof of Lemma 7.4, p.~68} \stcited{}, citing [Wi95, Cor.~1, 2, p.~482]. (2)
$H(M)$ is unramified at $q$, since $q\nmid Mp$ (good reduction of $X_1(Mp^r)$ at $q$). \stknown.
\emph{Proof} (seed sketch, made precise). Put $A=I_F/\cc_q(F)$. Since $\lambda(\fb)=\cc_q(F)$, the map $\lambda\bmod\cc_q(F)$ factors
through $R''$, and so through $\T(M)_\m\otimes I_F$ on the $T_\ell$, $\ell\nmid Npq$. So $\tr(\rho_F\otimes A)(\mathrm{Frob}_\ell)
=\tr(\rho_M\otimes A)(\mathrm{Frob}_\ell)$, where $\rho_M$ comes from $H(M)$. By Chebotarev and continuity the two traces agree on
$G_\Q$. Since $\rho_M$ is unramified at $q$, $\tr\rho_A(g\sigma)=\tr\rho_A(g)$ for $\sigma\in I_q$. By (H2), $A[G_\Q]\to M_2(A)$ is
surjective (Burnside modulo $\m_A$, then Nakayama). Since the trace form on $M_2(A)$ is perfect, $\rho_A(\sigma)=1$. \qed
\noindent The prime $q$ may be $\equiv-1\bmod p$ here: Prop.~\ref{prop:s06-galois} does not use (H0$^-$).

\subsubsection{Fixed-weight inputs and the case $q\equiv-1$}
\begin{theorem}[Kim--Ota]\label{thm:s06-KO}
Let $p\ge5$ and $f\in S_k(\Gamma_0(N))$ a newform with $N=N^+N^-$, and assume: (1) $(N,p)=1$; (2) $\bar\rho_f|_{G_{\Q(\sqrt{p^*})}}$
is absolutely irreducible; (3) $2\le k\le p-1$; (4) $(N^+,N^-)=1$; (5) $N^-$ squarefree; (6) if $q\mid N^-$ and $q\equiv\pm1\bmod p$,
then $\bar\rho_f$ is ramified at $q$. Then $\ord_\lambda\eta_f(N)=\ord_\lambda\eta_f(N^+,N^-)+\sum_{q\mid N^-}t_f(q)$, where the
congruence ideals are taken in the full and in the $N^-$-new Hecke algebras of $S_k(\Gamma_0(N))$, localised at $\m_f$.
\stcited{} [KO, Thm~1.1] \lit{kim-ota\_2019\_*}{Thm 1.1, p.~1}. At $k=2$ it is [PW, Thm~6.8]
\lit{pollack-weston\_2006\_*}{Thm 6.8, p.~21} with the definite self-pairing $\xi_f$, under PW's (CR), whose (2) is the same $\pm1$
condition \lit{pollack-weston\_2006\_*}{(CR), p.~2}.
\end{theorem}

\begin{corollary}[fixed-weight level lowering mod $\varpi^t$; replaces seed step (4)]\label{cor:s06-fixed}
Let $f$ be $q$-new of level $N=qM$ and weight $2\le k\le p-1$, with $p\nmid N$, (H2) for $\bar\rho_f$, $\m_f$ non-Eisenstein, and
(H0$^-$). Then $\cc_q(f)=(\varpi^{t_f(q)})$, and in particular $\cc_q(f)\subseteq(\varpi^{t_f(q)})$. \stpmc{}
\end{corollary}
\noindent\emph{Proof.} $\T_k(N)_{\m_f}$ is Gorenstein: [KO, Prop.~5.1(1)] \lit{kim-ota\_2019\_*}{Prop.~5.1, p.~18--19}, citing
[Dim09, Prop.~5.5]. We do not use [KO, Cor.~5.4] (Gorenstein property of the $N^-$-new quotient with no condition on $N^-$, in
conflict with [BKM] at trivial primes) beyond $q\not\equiv1$.
\provisional{D6: [KO, Cor.~5.4] restricted to $q\not\equiv1$; before: no statement (Cor.~5.4 was not cited in stages 05, 06).}
Thm~\ref{thm:s06-E}(iii) over the DVR $\cO$ (no condition on the new quotient) gives
$\eta_f(N)=\cc_q(f)\,\eta_f(M,q)$. Compare with Thm~\ref{thm:s06-KO} for $N^-=q$, $N^+=M$. \qed

\noindent So under (H0$^-$) step (4) of the seed does not need a new $R=\T$ argument (``passing to the quotient new at $\ell\,\|\,M$''). It
\emph{is} [KO, Thm~1.1] read through the fixed-weight Thm~E. Without (H0$^-$), the seed's direct $R=\T$ argument gives the inclusion
(Prop.~\ref{prop:s06-levM}). \provisional{F3: ``under (H0$^-$)'' and the pointer to the second proof added; before: step (4) ``is''
[KO, Thm~1.1], with no alternative.} The $R=\T$ theorems ([DFG, Thm~0.2 = 3.7]
\lit{diamond-flach-guo\_2004\_*}{Thm 0.2, p.~4}; [Dim09, Thm~6.6]
\lit{dimitrov\_2005\_iharas-lemma-for-hilbert-modular-varieties\_arxiv}{Thm 6.6, p.~23}; [KO, Thm~2.9, p.~5]) and
\textbf{Böckle's presentation} enter only inside the proof of [KO]. The latter is [KO, Thm~3.12 (Böckle), p.~12]: a complete
intersection presentation of the global deformation ring, citing [Böc07, Cor.~4.3, Thm~5.8] = G.~Böckle, \emph{Presentations of
universal deformation rings}, LMS Lecture Note Ser.\ 320 (2007), 24--58. That paper is not in \texttt{literature/}.

\begin{proposition}[the inclusion without (H0$^-$); the seed's step (4)]\label{prop:s06-levM}
Let $f$ be $q$-new of level $N=qM$ and weight $2\le k\le p-1$, with $p\nmid N$, (H0), (H2) for $\bar\rho_f$, $\m_f$ non-Eisenstein,
and (Lev$_M$). Then $\cc_q(f)\subseteq(\varpi^{t_f(q)})$. No condition on $q\bmod p$ beyond (H0) is used. \stpmc{}
\end{proposition}
\noindent\emph{Proof.} Put $t:=t_f(q)$. If $\bar\rho_f$ is ramified at $q$, then $t=0$ and there is nothing to prove.
(1) $\ord_\varpi\cc_q(f)$ is the largest $t'$ such that $\lambda_f\bmod\varpi^{t'}$ factors through $\T''_k$ (Def.~\ref{def:s06-ctam}:
$\cc_q(f)=\lambda_f(\fb_k)$).
(2) By Lemma~\ref{lem:s06-decomp}(b) at weight $k$, $\T''_k\cong\T_k(M)_\m$ with $U_q\mapsto\alpha_q$, the root of
$X^2-T_qX+q^{k-1}$ congruent to $U_q$. The two roots are $\equiv\varepsilon q^{(k-2)/2}$ and $\equiv\varepsilon q^{k/2}$, where
$a_q(f)=\varepsilon q^{(k-2)/2}$. They are distinct mod $\m$ because $q\not\equiv1$.
(3) $\rho_f\bmod\varpi^t$ is unramified at $q$ (definition of $t_f(q)$), low-weight crystalline at $p$, and of the same type as $f$
at the primes of $M$. So it is a point of the deformation ring that is minimal at $q$, ``new'' at the primes $\ell\,\|\,M$ where
$\bar\rho$ is unramified, and minimal elsewhere.
(4) By (Lev$_M$), [KO, Thm~3.14] \lit{kim-ota\_2019\_*}{Thm~3.14, p.~12; hypotheses: Assumption~2.1 and $P_{\bar\rho}\subseteq\Sigma$;
no condition $q\not\equiv\pm1$} identifies that ring with $\T_k(M)_\m$. So there is a map $\T_k(M)_\m\to\cO/\varpi^t$ that matches
$\lambda_f$ on the $T_\ell$.
(5) The image of $T_q$ is $\tr\rho_f(\Frob_q)=\varepsilon(q+1)q^{(k-2)/2}\bmod\varpi^t$. So the root $\alpha_q$ of (2) maps to
$\varepsilon q^{(k-2)/2}=a_q(f)$ (Hensel: the roots are distinct mod $\varpi$). Hence $\lambda_f\bmod\varpi^t$ factors through
$\T''_k$, and $\cc_q(f)\subseteq(\varpi^t)$ by (1). \qed

\noindent The prime $q$ does not enter the $R=\T$ theorem, because the deformation problem is unramified at $q$. This is the seed's
argument (``passing to the quotient that is new at the primes $\ell\,\|\,M$ where $\bar\rho$ is unramified and applying the $R=\T$
theorem there''). \textbf{Second reader required (D2):} step (4) must be checked against [KO, \S3, Thm~3.14] by a second reader before
the talk claims the case $q\equiv-1\bmod p$. Points to check: that the level-$M$ data of $f$ match Kim--Ota's level
$N^+_{\rho,\Sigma^+}N^-$ ($M$ squarefree at the primes where $\bar\rho$ is unramified, $f$ minimal at the others, $P_{\bar\rho}$ handled),
and that the map of (4) also matches $\lambda_f$ on the $U_\ell$, $\ell\mid M$ ([KO]'s Hecke algebra is the anemic one).
\provisional{F3(b)/D2: new Prop.~\ref{prop:s06-levM}, the five steps of audit issue 5 under (Lev$_M$), with a second-reader note
on step (4); before: no second proof; the seed's direct argument was declared unsupported.}

\begin{remark}[the case $q\equiv-1\bmod p$]\label{rem:s06-minus1}
The seed claims that step (4) covers $q\equiv-1\bmod p$. (a) [KO, Thm~1.1] as a black box excludes it: [KO, Thm~1.1(6)] and
[PW, (CR)(2)] both exclude $q\equiv\pm1$ with $\bar\rho$ unramified at $q$. So Cor.~\ref{cor:s06-fixed} needs (H0$^-$). (b) The
precise place in [KO] is the L-value step [KO, Prop.~2.15(1), p.~8] \lit{kim-ota\_2019\_*}{Prop.~2.15(1), p.~8; Prop.~5.6, p.~19;
Rem.~5.5, p.~19}. There, passing from the $\Sigma$-imprimitive to the primitive congruence ideal multiplies by
$L_q(\mathrm{ad}^0f,1)^{-1}=1-q^{-2}$ for $f$ special at $q$, and this is a $p$-unit iff $q\not\equiv\pm1$. The Selmer step
[KO, Thm~2.13, p.~7] needs only $q\not\equiv1$. This is a property of their method (they compare with the level where every lift at
$q$ is allowed), not of the statement $\cc_q(f)=(\varpi^{t_f(q)})$. (c) The seed's direct argument does not use [KO, Thm~1.1]; it is
Prop.~\ref{prop:s06-levM}, which needs only (Lev$_M$). For $q\equiv-1$ one still has $q+1\in\m$, which changes $\coker(\Psi^*\Psi)$
(Prop.~\ref{prop:s06-betti}). (d) Since (H0) already gives $q\not\equiv1$, the seed's exclusion ``$q\not\equiv\pm1$'' in T5 and C4(i)
is exactly (H0)$+$(H0$^-$). The other inputs at $q\equiv-1$: (Gor$_R$) has no condition at $q$, and (Gor$_{R'}$) needs only
(CR,$\{q\}$) (Remark~\ref{rem:s06-CR}). So Thm~E and Prop.~\ref{prop:s06-galois} hold, and Thm~\ref{thm:s06-defect} holds under
(Lev$_M$), in particular for $M=1$.
\provisional{F3(a),(d)/D2: ``seed step (4) is wrong as stated'' deleted; old item (c) (both $U_q$-signs congruent to old forms, so
$\T_\m$ sees only part of the local deformation) deleted as an obstruction; the ``Decision'' (theorem under (H0$^-$) only; equality
for $q\equiv-1$ \stopen, conjectured correction of size $v_\varpi(q+1)$ \stconj, HACS test via Risk R2) deleted; replaced by the
pointer to Prop.~levM.}
\end{remark}

\subsubsection{The defect formula for $D=q$}
\begin{theorem}[$\Lam$-adic defect formula; seed Thm~defect, corrected]\label{thm:s06-defect}
Assume (H0), (H1), (H2), (P) (with (P2) in its corrected form), and (H0$^-$) or (Lev$_M$). Then
\[(h_F)=\cc_q(F)=\tq_q(F)\qquad\text{as principal ideals of }I_F.\]
\stpmc{} in both cases (under (Lev$_M$): pending the second-reader check of Prop.~\ref{prop:s06-levM}, step (4)).
\provisional{F3(c)/D2: hypothesis ``(H0$^-$) or (Lev$_M$)''; before: (H0), (H0$^-$), (H1), (H2), (P).} Cited inputs: Hida freeness and control [Hi86b, Thm~3.1, Cor.~3.2]
\lit{hida\_1986\_iwasawa-modules-attached-to-congruences-of-cusp-forms}{Thm 3.1, Cor.~3.2 (3.3b), p.~250}; (Gor$_R$) and
(Gor$_{R'}$), with freeness of $H$ ([Hs] as above, resting on [Wi95], [MW86], [Hel, Cor.~8.11]
\lit{helm\_2004\_maps-between-modular-jacobians-and-jacobians-of-shimura-curves\_arxiv}{Cor.~8.11, p.~31}); the $\Lam$-adic JL
(\texttt{if:s01-targets}); and Thm~\ref{thm:s06-KO} at the point $P$ of (P) under (H0$^-$), or [KO, Thm~3.14] at $P$ under
(Lev$_M$). (Im) is not used.
\end{theorem}
\noindent\emph{Proof outline (corrected).}
(1) (Gor$_R$) and (Gor$_{R'}$) hold (as above), so $(h_F)=\cc_q(F)$ by Thm~\ref{thm:s06-E}.
(2) $\tq_q(F)\subseteq\cc_q(F)$ (Prop.~\ref{prop:s06-galois}).
(3) \emph{Control.} By Lemma~\ref{lem:s06-decomp}(b), $\T''\cong\T(M)_\m$ is $\Lam$-free. So by Lemma~\ref{lem:s06-D}(c), $\T'$ is
$\Lam$-free and $\fb$ is a summand. Hida's control for $\T$ and for $\T(M)_\m$ (which \emph{is} control for $\T''$, also at weight
$k$: $\T''_k\cong\T_k(M)_\m$), plus a rank count, gives
\provisional{F2: control for $\T''$ is control for $\T(M)_\m$ via the corrected Lemma decomp(b); before: ``Hida's control for $\T$
and for $\T(M)$ (hence for $\T''$)'', with $\T''$ the rank-two extension.}
$\T'/P_k\T'\cong\T'_k$ ($q$-new quotient at weight $k$) and $\fb\otimes\Lam/P_k=\fb_k$. Hence
$P(\cc_q(F))\cO_P=\cc_q(f)$. The seed states only ``$\subseteq$''; equality holds, but ``$\subseteq$'' suffices. For $k_P\ge3$
(case P1), $\T^{\ord}_k(Np)_\m\cong\T_k(N)_\m$ ($p$-stabilisation; $U_p$ is the unit root, by Hensel). For $k_P=2$ this is the
extra condition in (P2).
(4) Under (H0$^-$): Cor.~\ref{cor:s06-fixed}, $\cc_q(f)=(\varpi^{t_f(q)})$. Under (Lev$_M$): Prop.~\ref{prop:s06-levM},
$\cc_q(f)\subseteq(\varpi^{t_f(q)})$. The squeeze needs only the inclusion.
\provisional{F3(c): second route via Prop.~levM; before: ``Cor.~fixed: $\cc_q(f)=(\varpi^{t_f(q)})$. This needs (H0$^-$).''}
(5) \emph{Squeeze.} $P(\tq_q)\subseteq P(\cc_q)=\cc_q(f)\subseteq(\varpi^t)=P(\tq_q)$ (Lemma~\ref{lem:s06-tam}(b)), with $t<\infty$,
so all are equal. \provisional{F3(c): ``$=(\varpi^t)$'' weakened to ``$\subseteq(\varpi^t)$'', which is what step (4) gives under (Lev$_M$); before: $P(\tq_q)\subseteq P(\cc_q)=\cc_q(f)=(\varpi^t)=P(\tq_q)$.} Write
$\tq_q(F)=h_F\mathfrak d$ with $\mathfrak d\subseteq I_F$ (by (2)). Then $P(h_F)\ne0$ and $P(h_F)P(\mathfrak d)=P(h_F)\cO_P$, so
$P(\mathfrak d)=\cO_P$. Hence $\mathfrak d\not\subseteq\m_{I_F}$ and $\mathfrak d=I_F$. \qed

\begin{corollary}[all weights of the branch; seed Cor.~allweights]\label{cor:s06-allweights}
Under the hypotheses of Thm~\ref{thm:s06-defect}, let $P'$ be any arithmetic point of weight $k'\ge3$ with trivial character
(including $k'\ge p$), with $q$-new $f'=F_{P'}$. Then $\cc_q(f')=(\varpi^{t_{f'}(q)})$ and
$\ord_\varpi\eta_{f'}(N)-\ord_\varpi\eta_{f'}(M,q)=t_{f'}(q)$, for the ring-theoretic congruence ideals of $\T_{k'}(N)_\m$ and of its
$q$-new quotient. \stpmc{}
\end{corollary}
\noindent\emph{Proof.} Specialise Thm~\ref{thm:s06-defect} at $P'$. Use step (3) as an equality at $P'$, Lemma~\ref{lem:s06-tam}(b),
and Thm~\ref{thm:s06-E}(iii) over $\cO_{P'}$ with $\T_{k'}(N)_\m=\T/P'$, which is Gorenstein by Lemma~\ref{lem:s06-S}. For $k'\ge p$
this goes beyond [KO] ($k\le p-1$); it is the new content, together with the identification $(h_F)=\cc_q(F)$ on the new branch.

\begin{proposition}[Betti incarnation of $\cc_q$]\label{prop:s06-betti}
Let $l=p>k\ge4$ (the range of \texttt{if:s03-amalgam}), $q\not\equiv-1\bmod p$, and let $\m$ be non-Eisenstein for $\T(M)$ acting on
$H^1(\Gamma_0(M),V_k)$, with $H^1(\Gamma_0(M),V_k)_\m$ free of rank two over $\T(M)_\m$. Then:
\begin{enumerate}[label=(\alph*)]
\item $\coker(\Psi^*\Psi)_\m\cong\big(\T(M)_\m/(\delta)\big)^{\oplus2}$, where $\delta:=\det(\Psi^*\Psi)=(q+1)^2-q^{2-k}T_q^2$.
  Since $p\nmid2(q+1)$, $(\delta)=(T_q-\epsilon(q+1)q^{(k-2)/2})$ for the unique sign $\epsilon$ with this element in $\m$.
\item $\T_k(N)_{\m}/(\fa+\fb)\cong\T(M)_\m/(T_q-\epsilon(q+1)q^{(k-2)/2})$ (Ribet's level-raising module, with Ihara's lemma).
\item If $\rk_\cO\T'_{k,\m}=1$, then $\ord_\varpi\cc_q(f)=\tfrac12\,\mathrm{length}_\cO\coker(\Psi^*\Psi)_\m$. In general only
  $\ord_\varpi\cc_q(f)\le\tfrac12\mathrm{length}$ holds.
\end{enumerate}
(a) \sttoprove{} (linear algebra: $q+1$ is a unit, so $\coker$ of $\big(\begin{smallmatrix}q+1&*\\ *&q+1\end{smallmatrix}\big)$ on
$\T^2$ is $\T/\det$, and $\otimes_\T H^1_\m\cong\T^2$; it uses \emph{only} $\det$, so it is independent of the $T'_q/T''_q$ choice of
stage 03 D4). (b) \stcited{} for the existence of the congruences ([Ri84], not in \texttt{literature/}; [Dia89, Thm~4A--4C, Cor.~6.9]
\lit{diamond\_1989\_congruence-modules-associated-to-lambda-adic-forms}{Thm 4A p.~61, Cor.~6.9 p.~80}). The exact module structure is
\sttoprove. It is consistent with Lemma~\ref{lem:s06-decomp}(b): there $\T''_k\cong\T_k(M)_\m$ with $U_q\mapsto\alpha_q$, and
$T_q-\epsilon(q+1)q^{(k-2)/2}=(\alpha_q-\epsilon q^{(k-2)/2})(1-\epsilon q^{k/2}\alpha_q^{-1})$, where the second factor is
$\equiv1-q$ mod $\m$, a unit since $q\not\equiv1$. \provisional{F2: consistency of (b) with the corrected old quotient checked and
recorded; before: no remark.} (c) \stpmc: $\T'/\cc\cong\T/(\fa+\fb)$, and $\lambda'$ is an isomorphism when $\T'$ has rank one.
\end{proposition}
\noindent For $q\equiv-1$, (a) fails: the cokernel is no longer $(\T/\delta)^2$, and the two sign factors of $\delta$ both lie in
$\m$. For $k\ge p$ (T4 beyond $k=12$) \texttt{if:s03-amalgam} is outside its range $l>k$. One needs its version \emph{localised at the
non-Eisenstein $\m$ and after the ordinary projector}. There the end terms $H^0$ are Eisenstein and vanish (Risk R3). Stage 03 does not
export that version, so every use of the amalgam for $k\ge p$ (T4 at $k=28,44,60,124$) is \stexp.
\provisional{F19: uses of the amalgam for $k\ge p$ marked \stexp; before: no status for the $k\ge p$ uses.}

\begin{remark}[open cases; seed Remark open cases, O5]\label{rem:s06-open}
(i) $q\equiv-1\bmod p$ with $\bar\rho$ unramified at $q$, when (Lev$_M$) is not known, i.e.\ for $M>1$ (Remark~\ref{rem:s06-minus1},
Prop.~\ref{prop:s06-levM}). \provisional{F3: open case (i) restricted to $M>1$ without (Lev$_M$); before: all $q\equiv-1$ with
$\bar\rho$ unramified at $q$.} (ii) Branches whose only weight-two points
are $p$-new, or for which (P2)'s extra condition fails. The combinatorial half holds; the Galois half needs a version of [KO] at level
$Np$. (iii) Trivial primes, $q\equiv1\bmod p$ with $\bar\rho(\mathrm{Frob}_q)=\pm1$. Here the quaternionic Hecke algebra is
neither a complete intersection nor Gorenstein [BKM, Intro., Thm~1.2]
\lit{bockle-khare-manning\_2019\_wiles-defect-for-hecke-algebras-that-are-not-complete-intersections\_arxiv}{\S1, p.~2--3}
\stcited, so (Gor$_{R'}$) fails and $\cc$ is not principal (Thm~\ref{thm:s06-E}(iv)). By Thm~\ref{thm:s06-E}(iii) the identity
still holds at height-one primes. $q\equiv1$ non-trivial is excluded by (H0) through $p\nmid\varphi(N)$; [Hs] would allow it only
if $\bar\rho$ is ramified at $q$. (iv) Dropping (Im) (probably vacuous, R4). (v) Several primes in $D$.
\end{remark}

\begin{conjecture}[several ramified primes; seed \S10, last paragraph]\label{conj:s06-several}
Let $D=q_1\cdots q_r$ ($r$ odd), with each $q_i$ satisfying (H0), (H0$^-$) and $F$ $D$-new. Let
$\T=\T^{(0)}\twoheadrightarrow\T^{(1)}\twoheadrightarrow\cdots\twoheadrightarrow\T^{(r)}=\T^{D\text{-new}}\cong\T^{\ord}_{B_D}$ be the
$q_1\cdots q_i$-new quotients, and let $\cc^{(i)}_{q_i}(F)$ be the old/new congruence ideal of $\T^{(i-1)}\twoheadrightarrow\T^{(i)}$ at
$q_i$. Then $(h_F)=\prod_i\cc^{(i)}_{q_i}(F)=\prod_i\tq_{q_i}(F)$. \stconj{} If every $\T^{(i)}_\m$ is Gorenstein, the first
equality follows by iterating Thm~\ref{thm:s06-E} \stsketch. Whether the intermediate quotients are Gorenstein is \stopen{} (Hsieh
covers $i=0$ and $i=r$). The second equality would follow from Prop.~\ref{prop:s06-galois} at each step and [KO, Thm~1.1] with
$N^-=D$, by the squeeze of step (5) \stsketch. Note that $q_i$ is $q_1\cdots q_{i-1}$-new-ramified in the intermediate algebras,
which needs the relative version of Lemma~\ref{lem:s06-decomp} at each step.
\end{conjecture}

\subsection{Proof strategy}
\begin{itemize}
\item \emph{Commutative algebra} (Lemmas S, G, D, Thm~E) is complete above. It only remains to write it out, and the Lean track can
  start immediately.
\item \emph{Lemma~\ref{lem:s06-decomp}} is the one ``to prove'' input on the $\GL_2$ side that is not cited. Its delicate part is
  (b): Hensel factorisation of $X^2-T_qX+qS_q$ over $\T(M)_\m$ (roots distinct mod $\m$ since $q\not\equiv1$), so that $\T''\cong\T(M)_\m$
  is $\Lam$-free and $T_q$ lies in the image of $\T$. \provisional{F2: delicate part is now the Hensel factorisation; before: ``$T_q$
  lies in the image of $\T$ on the old space (Chebotarev argument) and $\T''$ is $\Lam$-free''.}
\item \emph{Galois half}: Burnside, Nakayama and Chebotarev, given the two cited freeness results.
\item \emph{Defect formula}: steps (1)--(5) above. The only fixed-weight arithmetic input is [KO, Thm~1.1] (under (H0$^-$)) or
  [KO, Thm~3.14] (under (Lev$_M$)), used at a single point. \provisional{F3: second fixed-weight input [KO, Thm~3.14]; before: only
  [KO, Thm~1.1].} Everything else is either commutative algebra or Hida control.
\item \emph{All weights}: specialisation, using control as an equality and Gorenstein base change.
\item \emph{Betti side}: linear algebra over $\T(M)_\m$, plus the cited level raising. Only $\det\Psi^*\Psi$ and the cokernel are used
  (stage 03, D4).
\end{itemize}

\subsection{Dependencies and interfaces}
\paragraph{Imported (restated verbatim from \texttt{SPLIT.md}).}
\begin{interface}[\texttt{if:s01-targets}]
Thm JLclassical (cited) and the $\Lam$-adic Hecke JL $\T^{\ord,D\text{-new}}\cong\T^\ord_B$ ([LV], [Hs]; cited).
\status{imported from stage 01}
\end{interface}
\begin{interface}[\texttt{if:s01-hecke}]
The Hecke operators $T_\ell$, $S_\ell$, $U_\ell$, $\langle d\rangle$, $U_p$ (with fixed normalisation $N_{U_p}$) and $\Pi_\ell$
($\ell\mid D$). All operators supported away from $q$ commute with right translation by $B_q^\times$. $T^*=T_{g^{-1}}$ for
$T=T_g$; $T_\ell$ is self-adjoint only at level $\Uzero{M}$.
\provisional{F15: exporter's wording; before: ``Self-adjointness for the pairing.''} \status{imported from stage 01}
\end{interface}
\begin{interface}[\texttt{if:s03-amalgam}]
The Mayer--Vietoris sequence for $\Gamma_0(M)[1/q]$, the formula for $\Psi^*\Psi$, and the definition of the Betti level-raising
module $\coker(\Psi^*\Psi)_\m$ (used by 06 as the Betti incarnation of $\cc_q$). \status{imported from stage 03}
\end{interface}
\begin{interface}[\texttt{if:s04-Mord}]
$\cM^\ord_B(U):=\varprojlim_se\,\Zp[X_B(UU_s)]$; $\Lam$-freeness and control (cited [Hi86b], [LV], [Hs]); and the Gorenstein property
of $\T^\ord_{B,\m}$ with $\cM^\ord_{B,\m}$ free of rank one (cited [Hs, Thm~4.5]) under the hypotheses ``$\m$ non-Eisenstein $+$
Hsieh's residual condition'' (stage 04 must state that condition explicitly). \status{imported from stage 04}
\end{interface}
\begin{interface}[\texttt{if:s04-etaB}]
$\eta^B_F$ equals the self-pairing of a generator of the $F$-eigenline (cited, Hs). \status{imported from stage 04}
\end{interface}
\begin{interface}[\texttt{if:s05-subquot} (optional)]
Prop subquot (also usable by 06). \status{imported from stage 05} Used only as a cross-check: with $H$ free of rank two over $\T$,
it gives $\coker(H[\fa]\to H/\fa H)\cong(\T'/\cc)^2$, the $\Lam$-adic analogue of the factor $2$ in Prop.~\ref{prop:s06-betti}(a).
(Stage 05 in turn imports our Lemma~\ref{lem:s06-D} through \texttt{if:s06-lemD}.) \provisional{F16: through the interface
\texttt{if:s06-lemD}; before: ``Stage 05 in turn uses our Lemma D as a citation.''}
\end{interface}

\paragraph{Exported (addition, not in \texttt{SPLIT.md}).}
\begin{interface}[\texttt{if:s06-lemD}]\label{if:s06-lemD}
Lemma~\ref{lem:s06-D} (a)--(c): new/old duality for a finite free $A$-algebra $S$ with ideals $\fa,\fb$ as there. (a) $\Ann_S\fa=\fb$,
$\Ann_S\fb=\fa$ (no Gorenstein hypothesis); (b) for $A$-Gorenstein $S$, $\fb\cong(S/\fa)^\vee$ and $\fa\cong(S/\fb)^\vee$; (c) if
moreover $S/\fb$ is $A$-free, $\fb$ is an $A$-summand and $S/\fa\cong\fb^\vee$ is $A$-free. Status \stproved. Imported by stage 05
(Prop.~subquot uses (a), Prop.~Xiso uses (b)).
\provisional{F16: new interface \texttt{if:s06-lemD}; before: Lemma~D was used by stage 05 as a plain citation, without an interface.}
\end{interface}
\begin{interface}[\texttt{if:s06-defect}]\label{if:s06-defect}
For the talk and Lean: Thm~\ref{thm:s06-E} with the explicit hypotheses (Gor$_S$), (Gor$_{S'}$) or the unimodularity condition;
Thm~\ref{thm:s06-defect} under (H0), (H1), (H2), (P) [with (P2) corrected], and (H0$^-$) or (Lev$_M$); Cor.~\ref{cor:s06-allweights}.
Status as stated there. The talk may claim the case $q\equiv-1\bmod p$ only after the second-reader check of
Prop.~\ref{prop:s06-levM}, step (4).
\provisional{F3/D2: ``(H0$^-$) or (Lev$_M$)'' and the talk caveat; before: under (H0), (H0$^-$), (H1), (H2), (P).}
\end{interface}

\paragraph{Checks against stages 01--05 (seam items).}
(S1) \emph{Factor $2$} (stage 01; stage 04 D5). The only level change in this stage is on the definite side, where T5 passes from
$\hat\cO^\times$ ($-1\in U$) to $\Uone{N}$, $N\ge3$ ($-1\notin U'$). There $\tr\circ\nu'=2\,\nu\circ\pi_*$; the hypothesis ``$-1\in U'$
or $-1\notin U$'' \emph{fails}. Since $p\ge5$, $2\in\Zp^\times$, so no ideal in this stage changes. Hsieh's weight is
$\#\Gamma_{Np^\alpha,x}=w_x$ (Hsieh (4.1), p.~29; stage 04); the factor $2$ is against the import's $|\Gamma_x|$. Integral Gram
matrices in T5 must record which of the two is used. \provisional{F10: Hsieh's weight is $w_x$; before: ``integral Gram matrices in T5
must record whether $w_x$ or $|\Gamma_x|$ is used (Hsieh uses $\#\Gamma$; stage 04 D4)''.}
(S2) \emph{Indefinite stage over $\Z$} (stage 02 D1, D3). This stage uses no indefinite tree stage, and everything is over $\Zp$,
$p\ge5$, where $H_1(\BG(U,q),\cdot)=H_1(\GU,\cdot)$. We never discharge $H_0(\GU,V)=V_{U^1_p}$ by a finite check.
(S3) \emph{Conventions} (stage 01). $U_p=[U\,\mathrm{diag}(p,1)\,U]$ on the definite side. The norm map is Hecke-equivariant only up
to adjoints, so $\eta^B_F$ computed from a Gram matrix uses the pairing $\langle Tx,y\rangle=\langle x,T^*y\rangle$. The congruence
ideal is unchanged: Hsieh's $B_N$ is Hecke-equivariant (the Atkin--Lehner twist turns adjoints into the operators themselves), so
$\eta^B_F$ is computed in $\T_B$ itself (stage 04 D4). \provisional{F18: $T\mapsto T^*$ is not $\Lam$-linear ($\langle d\rangle^*=
\langle d\rangle^{-1}$, semilinear for $[u]\mapsto[u^{-1}]$), so that argument is replaced; before: ``because $T\mapsto T^*$ is an
automorphism of $\T_B$ over $\Lam$ fixing $\ker\lambda'$ up to the Atkin--Lehner twist built into Hsieh's $B_N$''.} Adelic diamonds are inverse to classical ones (see \S\ref{ss:s06-setting}).
(S4) \emph{$\Psi^*\Psi$} (stage 03 D4): only $\det$ and $\coker$ are used.
(S5) \emph{$\Phicomb$ Eisenstein} (prop:s03-phieis): resolved; stage 03 proves it for every prime $q$ and every $M$ with
$\ell\nmid qM$, by an elementary argument. This stage does not use it. T4 ($q=3$) uses only the $\GL_2$ amalgam.
\provisional{F13: status ``resolved''; before: ``it is unresolved for $q\in\{2,3\}$''.}
(S6) \emph{Hsieh's residual condition} is ramification at $\ell\mid D$ with $\ell\equiv1$ only (stage 04 point (4); Remark~\ref{rem:s06-CR}).
(S7) \emph{Sign at $q$} (stage 05 D1). Stage 05 fixes $U_q\mapsto+\Pi_q$; Hsieh's opposite sign
\lit{hsieh\_2017\_*}{\S4.5, p.~34} is not imported (stage 05, Remark hsieh). The ring isomorphism $\T'\cong\T^\ord_B$ is determined
by the $T_\ell$, $\ell\nmid Np$, so no ideal depends on it. \provisional{F6/D3: Hsieh's sign is not imported; before: ``every statement
imported from [Hs] that involves $U_q$ must be re-signed''; the two conventions were said to name different images ($\pm\Pi_q$) of $U_q$
under the same ring isomorphism.} This stage does not use $\HypX$, so stage 05's (X3) factor-2 condition
and its $T_p$-at-$s=0$ and Rosati conventions do not enter.

\paragraph{Interface deviations.}
\begin{enumerate}[label=(D\arabic*)]
\item \texttt{if:s03-amalgam}: we use $\det(\Psi^*\Psi)=(q+1)^2-q^{2-k}T_q^2$ and $\coker(\Psi^*\Psi)_\m$ only. We also need a
  version localised at non-Eisenstein $\m$ and after $e$ for $k\ge p$ (T4), which stage 03 does not export (its MV needs $l>k$).
  Until it does, every use for $k\ge p$ is \stexp. \provisional{F19: status \stexp{} for $k\ge p$; before: no status.}
\item \texttt{if:s04-Mord}: ``Hsieh's residual condition'' is (CR,$\{q\}$) = absolutely irreducible $+$ $p$-distinguished $+$
  (ramified at $q$ if $q\equiv1$). Under (H0) it reduces to (CR). The $\Lam$-structures are matched by $[u]\mapsto[u]^{-1}$
  ($P_k\leftrightarrow P^{\rm ad}_k$, stage 04 D1).
\item \texttt{if:s04-etaB}: used through $\eta^B_F=\eta(R',\lambda')$ and the inclusion $(\eta^B_F)\supseteq(\eta_F)$ of [Hs, Rem.~7.8].
  The self-pairing is Hsieh's $B_N$, not the seed's $w_x\delta_{xy}$ (stage 04 D4). They agree up to units at $p\ge5$ on the
  ideal level only.
\item SPLIT lists no export for 06. We add \texttt{if:s06-lemD} (imported by stage 05) and \texttt{if:s06-defect}.
  \provisional{F16: \texttt{if:s06-lemD} added; before: ``We add \texttt{if:s06-defect}''.}
\end{enumerate}

\subsection{Lean, HACS, talk}
The tracks are independent. Each milestone refers only to statements that already exist above.

\paragraph{Mathematics track.}
\begin{enumerate}[label=M\arabic*.]
\item Write out Lemmas~\ref{lem:s06-S}, \ref{lem:s06-G}, \ref{lem:s06-D} and Thm~\ref{thm:s06-E}(i)--(iv) in full (proofs above).
\item Prove Lemma~\ref{lem:s06-decomp}, especially (b): the Hensel factorisation and $\T''\cong\T(M)_\m$ (hence $T_q$ in the image and
  $\Lam$-freeness of $\T''$). \provisional{F2: M2 follows the corrected (b); before: ``(b): $T_q$ in the image, and $\Lam$-freeness of
  $\T''$''.}
\item Write the control step (3) as an equality, including the $p$-stabilisation isomorphism, and settle (P2).
\item Prove Prop.~\ref{prop:s06-betti}(a), (c). Find [Ri84] for (b), or prove (b) from Ihara's lemma.
\item The case $q\equiv-1$: write out Prop.~\ref{prop:s06-levM}, have step (4) checked by a second reader against [KO, \S3,
  Thm~3.14], and determine for which $M>1$ (Lev$_M$) holds. \provisional{F3(d)/D2: M5 replaced; before: ``work out the expected
  correction term in Cor.~fixed via the Euler factor $1-q^{-2}$''.}
\item Conj.~\ref{conj:s06-several} for $r=3$: which intermediate quotients must be Gorenstein.
\end{enumerate}

\paragraph{Lean track.}
\begin{enumerate}[label=L\arabic*.]
\item \texttt{def IsGorensteinAlg (A S)}: $S$ finite free over $A$ and $\Hom_A(S,A)\cong S$ as $S$-modules; \texttt{congruenceIdeal}
  $\eta(S,\mu)$ as \texttt{Ideal A}.
\item Lemma S (without base change first), Lemma D(a), (b), Thm~E(i), (ii) over a commutative Noetherian local domain. These are pure
  module theory, Lean-friendly now.
\item Thm~E(iii) for $A$ a DVR (saturated $\Rightarrow$ summand over a PID; Mathlib has the structure theorem). Lemma G (Nakayama).
\item Algebraic core of Prop.~\ref{prop:s06-galois}: if $A$ is local, $A[G]\to M_2(A)$ is surjective, and $\tr\rho(g\sigma)=\tr\rho(g)$
  for all $g$, then $\rho(\sigma)=1$.
\item The squeeze, step (5): for a local domain $I$, ideals $\tq\subseteq(h)$ and a specialisation $\varphi:I\to\cO$ with
  $\varphi(\tq)=\varphi(h)\ne0$, conclude $\tq=(h)$.
\item Registered hypotheses (never axioms): \texttt{Hyp.HidaControl} (control at $P$, including the $p$-stabilisation
  isomorphism for $k_P\ge3$ and the no-$p$-new condition of (P2) for $k_P=2$), \texttt{Hyp.HsiehGorenstein} (Gor$_R$, Gor$_{R'}$),
  \texttt{Hyp.FreeRankTwo} ($H$ over $\T$), \texttt{Hyp.KimOta} (Thm~\ref{thm:s06-KO} at one point) or, under (Lev$_M$),
  \texttt{Hyp.KimOtaRT} ([KO, Thm~3.14] at one point), \texttt{Hyp.OldDecomp}
  (Lemma~\ref{lem:s06-decomp}, in the corrected form $\T''\cong\T(M)_\m$, $U_q\mapsto\alpha_q$). \provisional{F2/F3: \texttt{Hyp.OldDecomp}
  states the corrected (b); \texttt{Hyp.KimOtaRT} added for (Lev$_M$); before: \texttt{Hyp.OldDecomp} = Lemma decomp with
  $\T''\cong\T(M)_\m[X]/(X^2-T_qX+qS_q)$, no \texttt{Hyp.KimOtaRT}.} Then Thm~\ref{thm:s06-defect} becomes a Lean theorem modulo these.
\end{enumerate}

\paragraph{HACS / computation track.}
\begin{enumerate}[label=H\arabic*.]
\item \textbf{T4} ($q=3$, $p=17$, the branch through $g\equiv\Delta\bmod17$; hypotheses hold: $3\not\equiv\pm1\bmod17$, $17$ is not
  exceptional for $\Delta$). Compute $e\coker(\Psi^*\Psi)_\m$ at $k=12,28,44,60,124$ and check that it is $\cong(\T(1)_\m/\delta)^2$.
  Prediction: $\ord_{17}\cc_3=1,1,1,1,2$ (seed). For $k\ge17$ localise at $\m$ (including $T_{17}\notin\m$) and apply $e$ (D1).
  This tests Prop.~\ref{prop:s06-betti} and, where $t_{G_k}(3)$ is computed independently (Galois side, e.g.\ via
  $\ord_\varpi(a_3(G_k)^2-3^{k-2}\cdot16)$ on the old branch), Cor.~\ref{cor:s06-allweights} for $k\ge p$. The predictions at
  $k=28,44,60,124$ use the amalgam beyond its range and are \stexp{} (D1). \provisional{F19: $k\ge p$ predictions marked \stexp;
  before: no status.}
\item \textbf{T5} ($(q,k,p)\in\{(13,18,19),(5,18,23),(5,16,31),(13,12,31),(3,18,53),(3,18,61),(13,16,101)\}$; all satisfy
  $p\nmid q^2-1$ and $k\le p-1$). Compute the $q$-new Brandt module of $B_q$ with $\Sym^{k-2}$ at $\m$, its Gram matrix with Hsieh's
  weights (Hsieh's weight is $w_x$; the factor $2$ is against the import's $|\Gamma_x|$, S1) \provisional{F10: before: ``($\#\Gamma_x$;
  record the factor $2$, S1)''.}, $\eta^B$ from the Gram determinant ([Hs, Rem.~7.3, p.~67]:
  $C(\lambda)^d=\det\langle e_i,e_j\rangle$), and $\eta_f(q)$ and $\cc_q(f)$ from HACS. Check $\eta_f(q)=\cc_q(f)\eta^B$. Note: at
  $k\le p-1$ this is \emph{predicted} by Thm~E (proved) and [KO]; it validates the definite module and the identification of $\eta^B$
  with the Gram self-pairing. It is genuinely new only at $k+(p-1)\ge p$ (Cor.~\ref{cor:s06-allweights}).
\item \textbf{[validation] $q\equiv-1$}: repeat T5 for triples with $p\mid q+1$ (e.g.\ $q=13$, $p=7$, $3\le k\le6$; $q=19$ or $q=29$ with $p=5$, $k=4$;
  all satisfy $p\nmid\varphi(q)$) to measure $\ord\cc_q(f)-t_f(q)$. This validates Prop.~\ref{prop:s06-levM} (here $M=1$). Predicted
  difference: $0$. The Hecke algebra used must contain $U_q$. \provisional{F3(d): H3 is a validation with predicted difference $0$
  and $U_q$ in the Hecke algebra; before: ``New, $q\equiv-1$'', measuring a conjectured correction (Remark minus1).}
\item C1, C3, C4(i) are done; C4(i)'s restriction $q\not\equiv\pm1$ is exactly (H0)$+$(H0$^-$).
\end{enumerate}

\paragraph{Talk track.}
\begin{enumerate}[label=K\arabic*.]
\item ``The defect is an old/new congruence ideal'': $\eta_F=\cc_q(F)\cdot\eta^B_F$, a three-line proof (Thm~E). The two Gorenstein
  hypotheses are where the arithmetic hides.
\item ``\dots and it equals a Tamagawa ideal'': the squeeze with one Kim--Ota point. The numbers $17$ ($q=3$), $29,131$ ($q=5$) from
  HACS (C1, C3).
\item The caveat slide: $q\equiv-1\bmod p$ (covered for $M=1$ by Prop.~\ref{prop:s06-levM}, only after the second-reader check)
  and trivial primes (BKM), and why they are different. \provisional{F3/D2: caveat for $q\equiv-1$ updated; before: ``$q\equiv-1\bmod p$
  and trivial primes (BKM), and why they are different''.}
\end{enumerate}

\subsection{Risks}
\begin{enumerate}[label=R\arabic*.]
\item \emph{Seed step (4) at $q\equiv-1$} rests on [KO, Thm~3.14] through (Lev$_M$) (Prop.~\ref{prop:s06-levM}). Step (4) of that
  proof has not been checked by a second reader. Mitigation: (H0$^-$) as the alternative hypothesis, the second-reader check, and
  the HACS test H3. \provisional{F3/D2: before: ``Seed step (4) at $q\equiv-1$ is unsupported. Mitigation: (H0$^-$), and the HACS
  test H3.''}
\item (Deleted.) \provisional{F3(d): risk ``expected correction at $q\equiv-1$'' deleted; before: a guessed correction of size
  $v(q+1)$ from the Euler factor, status at most \stopen.}
\item \emph{T4 beyond $k=p$} is outside the stage-03 range $l>k$ (and outside the perfectness range $p>k-2$ of the $V_k$-pairing,
  stage 01 risk). It needs the localised and ordinary version of the amalgam computation (D1).
\item \emph{(Im) is probably superfluous.} Its source in the seed is unclear (perhaps an earlier version of [DFG]'s condition $\lambda\notin
  S_f$). Keeping it is harmless; removing it requires checking every cited theorem (done here for [KO], [Hs], [Wi95]).
\item \emph{(P2) for $k_0\equiv2$}: without the no-$p$-new condition the comparison with [KO] at level $N$ fails. Branches through
  weight-two points with $a_p\equiv\pm1$ are at risk.
\item \emph{$\Lam$-structures and pairings}: diamond inversion and Hsieh's $B_N$ versus $w_x\delta_{xy}$ affect only units at
  $p\ge5$. Misapplied, they would send $P_k$ to $P_{4-k}$ (S3, D2).
\item \emph{Literature gaps}: [Ri84], [MW86], [Böc07], [Hi86a] are missing; see \texttt{missing/stage-06.md}.
  [Hi86b] is a scan with poor OCR (theorem numbers were checked, statements only partially).
\end{enumerate}
